\documentclass[11pt,a4paper]{article}
\usepackage[margin=2.5cm]{geometry}
\usepackage{amsmath,amssymb}
\usepackage{booktabs}
\usepackage{graphicx}
\usepackage{xcolor}
\usepackage{hyperref}
\usepackage{parskip}
\usepackage{array}
\usepackage{multirow}

\definecolor{warnred}{RGB}{180,0,0}
\definecolor{okgreen}{RGB}{0,130,0}
\definecolor{warnorange}{RGB}{200,100,0}

\hypersetup{colorlinks=true,linkcolor=blue,urlcolor=blue}

\title{\textbf{Calibration Scan Results: Rank-1 Covariance Inflation Amplitude}\\[0.4em]
       \large \texttt{coh\_amplitude\_scan} — Run \texttt{2026-05-22T16:23:46}}
\author{destacking project}
\date{2026 May 22--23 \quad (7\,h\,45\,m wall time)}

\begin{document}
\maketitle

\begin{abstract}
We ran a calibration scan sweeping the rank-1 covariance inflation amplitude
(\texttt{coh\_amplitude} $\varepsilon \in \{0.0,\,0.1,\,0.2,\,0.3,\,0.4,\,0.6\}$)
with 10 independent realisations of 5000-cluster DES Y3-like simulations per amplitude.
The target diagnostic is the calibration ratio
$\sigma_{\rm runs}/\langle\sigma_{\rm post}\rangle$, where unity indicates that
posterior error bars correctly reflect run-to-run scatter.
\textbf{No amplitude achieved a calibration ratio near 1.}
The base case ($\varepsilon=0$) already gives a ratio of $1.97$ (posteriors
$\sim$2$\times$ too narrow).  At $\varepsilon \ge 0.2$, a systematic spurious
posterior mode at $\hat\gamma \approx +1.4$--$1.6$ progressively contaminates
results, rendering all calibration measurements at those amplitudes
unreliable.  By $\varepsilon = 0.4$--$0.6$ every realisation is dominated
by this mode, many with degenerate single-point posteriors ($\sigma=0$).
The spurious mode must be identified and eliminated before the amplitude scan
can be used for calibration.
\end{abstract}

\tableofcontents
\bigskip

% ============================================================
\section{Scan configuration}
% ============================================================

\begin{tabular}{ll}
\toprule
Parameter & Value \\
\midrule
\texttt{N\_CLUSTERS}      & 5000 per realisation \\
\texttt{N\_REAL\_PER\_AMP} & 10 \\
\texttt{N\_WORKERS}        & 15 \\
\texttt{AMPLITUDES}        & 0.0, 0.1, 0.2, 0.3, 0.4, 0.6 \\
True $\gamma$              & 0.2 \\
\texttt{DELTA\_M\_REF}     & 0.1 mag \\
\texttt{N\_LIVE} (nautilus) & 400 \\
Run directory              & \texttt{coh\_tuning\_outputs/run\_2026-05-22T16-23-46} \\
Total wall time            & 7 h 45 m \\
\bottomrule
\end{tabular}

% ============================================================
\section{Aggregate results}
% ============================================================

Table~\ref{tab:aggregate} shows the per-amplitude summary statistics produced
by \texttt{coh\_amplitude\_scan.py}.  The calibration ratio
$r = \sigma_{\rm runs}/\langle\sigma_{\rm post}\rangle$ is the primary
diagnostic; $r=1$ is the target.

\begin{table}[h]
\centering
\caption{Aggregate calibration results. $N=10$ realisations per amplitude.
  $\bar{\hat\gamma}$ = mean recovered $\gamma$;
  $\sigma_{\rm runs}$ = std of $\hat\gamma$ across realisations;
  $\langle\sigma_{\rm post}\rangle$ = mean posterior half-width;
  $r$ = calibration ratio;
  $\bar b$ = mean bias $\bar{\hat\gamma} - \gamma_{\rm true}$.}
\label{tab:aggregate}
\begin{tabular}{r r r r r r r}
\toprule
$\varepsilon$ & $N$ & $\bar{\hat\gamma}$ & $\sigma_{\rm runs}$ &
  $\langle\sigma_{\rm post}\rangle$ & $r$ & $\bar{b}$ \\
\midrule
\textcolor{okgreen}{0.000} & 10 & $-0.056$ & 0.292 & 0.148 & \textcolor{okgreen}{1.97} & $-0.256$ \\
\textcolor{warnorange}{0.100} & 10 & $+0.135$ & 0.487 & 0.141 & \textcolor{warnorange}{3.46} & $-0.065$ \\
\textcolor{warnred}{0.200} & 10 & $+0.636$ & 0.665 & 0.110 & \textcolor{warnred}{6.02} & $+0.437$ \\
\textcolor{warnred}{0.300} & 10 & $+1.004$ & 0.711 & 0.070 & \textcolor{warnred}{10.21} & $+0.804$ \\
\textcolor{warnred}{0.400} & 10 & $+1.204$ & 0.656 & 0.021 & \textcolor{warnred}{31.58} & $+1.004$ \\
\textcolor{warnred}{0.600} & 10 & $+1.359$ & 0.952 & 0.017 & \textcolor{warnred}{54.84} & $+1.159$ \\
\bottomrule
\end{tabular}
\end{table}

The three diagnostic plots produced by the scan are reproduced in
Figures~\ref{fig:calibration}--\ref{fig:bias}.

\begin{figure}[h]
\centering
\includegraphics[width=0.85\textwidth]{coh_amplitude_scan_calibration.png}
\caption{Calibration ratio $r = \sigma_{\rm runs}/\langle\sigma_{\rm post}\rangle$
vs.\ \texttt{coh\_amplitude}.
The target $r=1$ (dashed) is not achieved at any amplitude.
$r$ increases monotonically from $1.97$ at $\varepsilon=0$ to $54.8$ at
$\varepsilon=0.6$, driven by the collapse of $\langle\sigma_{\rm post}\rangle$
as the spurious mode dominates.}
\label{fig:calibration}
\end{figure}

\begin{figure}[h]
\centering
\includegraphics[width=0.85\textwidth]{coh_amplitude_scan_widths.png}
\caption{Empirical run-to-run scatter $\sigma_{\rm runs}$ (red) and mean
posterior width $\langle\sigma_{\rm post}\rangle$ (blue) vs.\ amplitude.
$\langle\sigma_{\rm post}\rangle$ collapses from $\sim$0.14 at $\varepsilon \le 0.1$
to $\sim$0.017 at $\varepsilon=0.6$ as degenerate posteriors dominate.}
\label{fig:widths}
\end{figure}

\begin{figure}[h]
\centering
\includegraphics[width=0.85\textwidth]{coh_amplitude_scan_bias.png}
\caption{Mean recovered $\gamma$ bias $\bar{b} = \bar{\hat\gamma} - \gamma_{\rm true}$
vs.\ amplitude.  Bias is already $-0.26$ at $\varepsilon=0$ and rises steeply
above $\varepsilon=0.1$ as the spurious $\hat\gamma \approx +1.4$ mode
dominates the sample mean.}
\label{fig:bias}
\end{figure}

% ============================================================
\section{Per-realisation data}
% ============================================================

Table~\ref{tab:perreal} lists every realisation's recovered $\hat\gamma$ and
posterior half-width $\sigma_{\rm post}$.  Entries highlighted in red are
identified as affected by the spurious mode
($|\hat\gamma - \gamma_{\rm true}| > 0.5$ and/or $\sigma_{\rm post} < 0.05$).

\begin{table}[h]
\centering
\caption{Per-realisation results. $\gamma_{\rm true} = 0.2$.
  Red = affected by spurious mode.}
\label{tab:perreal}
\small
\begin{tabular}{r l r@{\,}l r}
\toprule
$\varepsilon$ & Real.\ & \multicolumn{2}{c}{$\hat\gamma\;[68\%\,\rm CI]$} & $\sigma_{\rm post}$ \\
\midrule
\multirow{10}{*}{0.000}
 & 000 & $-0.496$ & $[-0.596,\,-0.368]$ & 0.114 \\
 & 001 & $+0.021$ & $[-0.134,\,+0.190]$ & 0.162 \\
 & 002 & $+0.385$ & $[+0.217,\,+0.568]$ & 0.175 \\
 & 003 & $+0.251$ & $[+0.060,\,+0.415]$ & 0.178 \\
 & 004 & $+0.107$ & $[-0.044,\,+0.259]$ & 0.151 \\
 & 005 & $-0.245$ & $[-0.394,\,-0.100]$ & 0.147 \\
 & 006 & $+0.163$ & $[+0.025,\,+0.301]$ & 0.138 \\
 & 007 & $-0.056$ & $[-0.200,\,+0.097]$ & 0.148 \\
 & 008 & $-0.298$ & $[-0.436,\,-0.129]$ & 0.154 \\
 & 009 & $-0.390$ & $[-0.505,\,-0.271]$ & 0.117 \\
\midrule
\multirow{10}{*}{0.100}
 & 000 & $+0.638$ & $[+0.489,\,+0.813]$ & 0.162 \\
 & 001 & $+0.372$ & $[+0.239,\,+0.512]$ & 0.137 \\
 & 002 & $+0.379$ & $[+0.212,\,+0.586]$ & 0.187 \\
 & 003 & $+0.104$ & $[-0.059,\,+0.271]$ & 0.165 \\
 & 004 & $+1.062$ & $[+0.940,\,+1.163]$ & 0.111 \\
 & 005 & $-0.572$ & $[-0.631,\,-0.500]$ & 0.065 \\
 & 006 & $-0.260$ & $[-0.435,\,-0.091]$ & 0.172 \\
 & 007 & $+0.020$ & $[-0.142,\,+0.177]$ & 0.159 \\
 & 008 & $-0.098$ & $[-0.216,\,+0.025]$ & 0.121 \\
 & 009 & $-0.290$ & $[-0.392,\,-0.135]$ & 0.129 \\
\midrule
\multirow{10}{*}{0.200}
 & 000 & $-0.221$ & $[-0.345,\,-0.084]$ & 0.131 \\
 & 001 & $+0.614$ & $[+0.411,\,+0.751]$ & 0.171 \\
 & 002 & $+0.326$ & $[+0.187,\,+0.449]$ & 0.131 \\
 & 003 & $+0.432$ & $[+0.314,\,+0.567]$ & 0.127 \\
 & 004 & $+0.415$ & $[+0.262,\,+0.574]$ & 0.156 \\
 & 005 & $-0.053$ & $[-0.210,\,+0.105]$ & 0.157 \\
 & \textcolor{warnred}{006} & \textcolor{warnred}{$+1.436$} & \textcolor{warnred}{$[+1.419,\,+1.453]$} & \textcolor{warnred}{0.017} \\
 & \textcolor{warnred}{007} & \textcolor{warnred}{$+1.516$} & \textcolor{warnred}{$[+1.516,\,+1.560]$} & \textcolor{warnred}{0.022} \\
 & 008 & $+0.254$ & $[+0.083,\,+0.416]$ & 0.166 \\
 & \textcolor{warnred}{009} & \textcolor{warnred}{$+1.648$} & \textcolor{warnred}{$[+1.647,\,+1.702]$} & \textcolor{warnred}{0.027} \\
\midrule
\multirow{10}{*}{0.300}
 & \textcolor{warnred}{000} & \textcolor{warnred}{$+1.162$} & \textcolor{warnred}{$[+1.060,\,+1.256]$} & \textcolor{warnred}{0.098} \\
 & 001 & $+0.406$ & $[+0.276,\,+0.515]$ & 0.119 \\
 & 002 & $-0.802$ & $[-0.833,\,-0.619]$ & 0.107 \\
 & \textcolor{warnred}{003} & \textcolor{warnred}{$+1.307$} & \textcolor{warnred}{$[+1.261,\,+1.340]$} & \textcolor{warnred}{0.039} \\
 & \textcolor{warnred}{004} & \textcolor{warnred}{$+1.319$} & \textcolor{warnred}{$[+1.118,\,+1.388]$} & \textcolor{warnred}{0.135} \\
 & \textcolor{warnred}{005} & \textcolor{warnred}{$+1.560$} & \textcolor{warnred}{$[+1.560,\,+1.563]$} & \textcolor{warnred}{0.002} \\
 & \textcolor{warnred}{006} & \textcolor{warnred}{$+1.357$} & \textcolor{warnred}{$[+1.313,\,+1.402]$} & \textcolor{warnred}{0.045} \\
 & \textcolor{warnred}{007} & \textcolor{warnred}{$+0.981$} & \textcolor{warnred}{$[+0.912,\,+1.049]$} & \textcolor{warnred}{0.068} \\
 & \textcolor{warnred}{008} & \textcolor{warnred}{$+1.344$} & \textcolor{warnred}{$[+1.292,\,+1.393]$} & \textcolor{warnred}{0.051} \\
 & \textcolor{warnred}{009} & \textcolor{warnred}{$+1.410$} & \textcolor{warnred}{$[+1.377,\,+1.441]$} & \textcolor{warnred}{0.032} \\
\midrule
\multirow{10}{*}{0.400}
 & \textcolor{warnred}{000} & \textcolor{warnred}{$+1.545$} & \textcolor{warnred}{$[+1.545,\,+1.545]$} & \textcolor{warnred}{0.000} \\
 & \textcolor{warnred}{001} & \textcolor{warnred}{$+1.390$} & \textcolor{warnred}{$[+1.356,\,+1.415]$} & \textcolor{warnred}{0.030} \\
 & \textcolor{warnred}{002} & \textcolor{warnred}{$+1.432$} & \textcolor{warnred}{$[+1.423,\,+1.440]$} & \textcolor{warnred}{0.008} \\
 & \textcolor{warnred}{003} & \textcolor{warnred}{$+1.384$} & \textcolor{warnred}{$[+1.350,\,+1.419]$} & \textcolor{warnred}{0.035} \\
 & \textcolor{warnred}{004} & \textcolor{warnred}{$+1.476$} & \textcolor{warnred}{$[+1.476,\,+1.476]$} & \textcolor{warnred}{0.000} \\
 & \textcolor{warnred}{005} & \textcolor{warnred}{$+0.929$} & \textcolor{warnred}{$[+0.929,\,+0.929]$} & \textcolor{warnred}{0.000} \\
 & \textcolor{warnred}{006} & \textcolor{warnred}{$+1.212$} & \textcolor{warnred}{$[+1.091,\,+1.291]$} & \textcolor{warnred}{0.100} \\
 & \textcolor{warnred}{007} & \textcolor{warnred}{$-0.543$} & \textcolor{warnred}{$[-0.543,\,-0.543]$} & \textcolor{warnred}{0.000} \\
 & \textcolor{warnred}{008} & \textcolor{warnred}{$+1.845$} & \textcolor{warnred}{$[+1.845,\,+1.845]$} & \textcolor{warnred}{0.000} \\
 & \textcolor{warnred}{009} & \textcolor{warnred}{$+1.367$} & \textcolor{warnred}{$[+1.331,\,+1.402]$} & \textcolor{warnred}{0.035} \\
\midrule
\multirow{10}{*}{0.600}
 & \textcolor{warnred}{000} & \textcolor{warnred}{$+1.605$} & \textcolor{warnred}{$[+1.605,\,+1.677]$} & \textcolor{warnred}{0.036} \\
 & \textcolor{warnred}{001} & \textcolor{warnred}{$+1.529$} & \textcolor{warnred}{$[+1.529,\,+1.529]$} & \textcolor{warnred}{0.000} \\
 & \textcolor{warnred}{002} & \textcolor{warnred}{$+1.501$} & \textcolor{warnred}{$[+1.414,\,+1.501]$} & \textcolor{warnred}{0.044} \\
 & \textcolor{warnred}{003} & \textcolor{warnred}{$+2.188$} & \textcolor{warnred}{$[+2.188,\,+2.188]$} & \textcolor{warnred}{0.000} \\
 & \textcolor{warnred}{004} & \textcolor{warnred}{$+1.665$} & \textcolor{warnred}{$[+1.665,\,+1.665]$} & \textcolor{warnred}{0.000} \\
 & \textcolor{warnred}{005} & \textcolor{warnred}{$+1.566$} & \textcolor{warnred}{$[+1.566,\,+1.566]$} & \textcolor{warnred}{0.000} \\
 & \textcolor{warnred}{006} & \textcolor{warnred}{$+1.692$} & \textcolor{warnred}{$[+1.692,\,+1.692]$} & \textcolor{warnred}{0.000} \\
 & \textcolor{warnred}{007} & \textcolor{warnred}{$+1.619$} & \textcolor{warnred}{$[+1.619,\,+1.619]$} & \textcolor{warnred}{0.000} \\
 & \textcolor{warnred}{008} & \textcolor{warnred}{$-1.292$} & \textcolor{warnred}{$[-1.480,\,-1.292]$} & \textcolor{warnred}{0.094} \\
 & \textcolor{warnred}{009} & \textcolor{warnred}{$+1.515$} & \textcolor{warnred}{$[+1.515,\,+1.515]$} & \textcolor{warnred}{0.000} \\
\bottomrule
\end{tabular}
\end{table}

% ============================================================
\section{Central finding: spurious posterior mode at $\hat\gamma \approx +1.4$}
% ============================================================

\subsection{Phenomenology}

Starting at $\varepsilon = 0.2$, a fraction of realisations produce posteriors
concentrated near $\hat\gamma \approx +1.4$--$1.6$ with anomalously narrow
widths ($\sigma_{\rm post} \approx 0.002$--$0.044$, compared to the expected
$\sim$0.13--0.18 from $\varepsilon=0$ realisations).
The affected fraction grows with amplitude:

\begin{center}
\begin{tabular}{r c l}
\toprule
$\varepsilon$ & Affected / total & Notes \\
\midrule
0.000 & 0/10 & — \\
0.100 & 0/10$^*$ & $^*$scatter increased but no degenerate posteriors \\
0.200 & 3/10 & $\hat\gamma = +1.44,\,+1.52,\,+1.65$ \\
0.300 & 8/10 & systematic; only 2 realisations near truth \\
0.400 & 10/10 & all affected; 5 with $\sigma=0$ \\
0.600 & 10/10 & all affected; 7 with $\sigma=0$; one at $\hat\gamma=+2.19$ \\
\bottomrule
\end{tabular}
\end{center}

\subsection{Degenerate posteriors}

The $\sigma=0$ cases indicate that the nautilus sampler converged to a single
point in parameter space.  This is symptomatic of a very sharp likelihood spike
— the sampler finds the mode quickly and the neural-network surrogate declares
convergence before the peak is resolved.  These are \emph{not} legitimate
tight posteriors.

\subsection{Likely cause}

With rank-1-inflated covariance
$\mathbf{C}_i^{\rm eff} = \mathbf{C}_i + \varepsilon^2 \mathbf{m}_i\mathbf{m}_i^\top$,
the effective log-likelihood is
\begin{equation}
  \ln\mathcal{L}_i \propto
  -\tfrac12 \left[
    \frac{(\mathbf{y}_i\cdot\hat{\mathbf{w}}_i)^2}{1+\|\mathbf{w}_i\|^2}
    + \|\mathbf{y}_{i,\perp}\|^2
  \right]
  - \tfrac12 \ln(1 + \|\mathbf{w}_i\|^2),
\end{equation}
where $\hat{\mathbf{w}}_i = \mathbf{L}_i^{-1}\varepsilon\mathbf{m}_i /
\|\mathbf{L}_i^{-1}\varepsilon\mathbf{m}_i\|$ and
$\mathbf{y}_{i,\perp}$ is the residual component perpendicular to
$\mathbf{m}_i$.  At large $\varepsilon$, the parallel component is
heavily suppressed ($\to 0$), and the log-determinant penalty grows as
$\tfrac12\ln(\varepsilon^2\|\mathbf{L}_i^{-1}\mathbf{m}_i\|^2)$.

There exists a parameter regime (near $\gamma\approx+1.4$) where the model
predictions $\mathbf{m}_i$ happen to be nearly perpendicular to the data
residuals under the base covariance, while the log-det penalty is relatively
small.  This creates a secondary likelihood mode that becomes the global
maximum for $\varepsilon \gtrsim 0.2$.

This is a structural consequence of the rank-1 parameterisation: the assumed
coherent variance direction (proportional to $\mathbf{m}$) does not match
the actual coherent variance in the simulation (triaxiality $+$ miscentering),
creating a self-consistency problem at large $\varepsilon$.

% ============================================================
\section{Base-case undercoverage ($\varepsilon = 0$)}
% ============================================================

Even without rank-1 inflation, the calibration ratio is $r = 1.97$: the
posterior widths are roughly $2\times$ too narrow.  The run-to-run scatter of
$\hat\gamma$ ($\sigma_{\rm runs} = 0.29$) is larger than the typical posterior
half-width ($\langle\sigma_{\rm post}\rangle = 0.15$), indicating that the
per-realisation posteriors do not encompass the true variability of the
estimator.

This undercoverage at $\varepsilon=0$ is a separate problem from the spurious
mode, and needs to be addressed regardless of the rank-1 inflation strategy.
Possible sources include: (a) miscentering and triaxiality scatter that is not
captured by the spherical NFW likelihood; (b) sample noise from using only
$N=10$ realisations.

\section{Reliability of calibration ratios}

\begin{center}
\begin{tabular}{r c l}
\toprule
$\varepsilon$ & Calibration ratio & Reliability \\
\midrule
0.000 & 1.97 & \textcolor{okgreen}{Reliable — no spurious mode} \\
0.100 & 3.46 & \textcolor{warnorange}{Marginal — scatter inflated by amp=0.1 mode onset} \\
0.200 & 6.02 & \textcolor{warnred}{Unreliable — 3/10 realisations affected} \\
0.300 & 10.21 & \textcolor{warnred}{Unreliable — 8/10 realisations affected} \\
0.400 & 31.58 & \textcolor{warnred}{Unreliable — all realisations affected} \\
0.600 & 54.84 & \textcolor{warnred}{Unreliable — all realisations affected} \\
\bottomrule
\end{tabular}
\end{center}

% ============================================================
\section{Recommendations}
% ============================================================

\begin{enumerate}

\item \textbf{Diagnose the $\hat\gamma \approx +1.4$ mode.}
  Compute the log-likelihood landscape as a function of $\gamma$ alone
  (marginalising or fixing other parameters) for a single realisation at
  $\varepsilon = 0.2$, comparing $\varepsilon=0$ and $\varepsilon=0.2$ profiles.
  Identify whether the mode arises from the log-det penalty, the parallel
  suppression term, or a combination.

\item \textbf{Re-examine the rank-1 parameterisation.}
  The assumption $\mathbf{u}_i = \varepsilon\,\mathbf{m}_i(\theta)$ ties the
  coherent variance direction to the current model prediction.  A fixed,
  amplitude-independent coherent direction (e.g.\ the mean profile at a
  reference parameter value) may avoid the parameter-dependent mode.

\item \textbf{Increase $N_{\rm real}$ for $\varepsilon=0$.}
  With only 10 realisations, $\sigma_{\rm runs}$ has $\sim$23\% uncertainty.
  Running $\sim$50 realisations at $\varepsilon=0$ would give a reliable
  baseline calibration ratio and confirm whether the undercoverage ($r\approx2$)
  is real or a small-sample artefact.

\item \textbf{Do not use $\varepsilon \ge 0.2$ until the mode is fixed.}
  The calibration ratios at $\varepsilon = 0.2$--$0.6$ are dominated by the
  spurious mode and have no physical meaning as calibration diagnostics.

\end{enumerate}

% ============================================================
\section*{Appendix: Run notes}
% ============================================================

\begin{itemize}
  \item A harmless \texttt{OSError: [Errno 9] Bad file descriptor} appears at
    the end of the log from the \texttt{multiprocessing.Pool} destructor — not
    a scan failure.
  \item Total failures (realisations that raised exceptions and were retried):
    $\varepsilon=0.3$: 2 failures; $\varepsilon=0.4$: 0; $\varepsilon=0.6$: 4 failures.
  \item The previous scan run (\texttt{run\_2026-05-21T10-17-02}) produced
    entirely invalid results for $\varepsilon>0$ due to catastrophic cancellation
    in the Sherman--Morrison quadratic form (described in
    \texttt{2026-05-22\_coh\_amplitude\_scan\_fix.pdf}).
    The present scan uses the numerically stable fix committed to
    \texttt{destacking/fitting.py}.
\end{itemize}

\end{document}
